\documentclass[11pt]{article}
\usepackage[T1]{fontenc}
\usepackage{lmodern}
\usepackage[margin=1in]{geometry}
\usepackage{amsmath,amssymb,amsthm,booktabs,array}
\usepackage[hidelinks]{hyperref}
\newtheorem{thm}{Theorem}[section]
\newtheorem{lem}[thm]{Lemma}
\newtheorem{cor}[thm]{Corollary}
\newcommand{\gt}{\gamma_{\mathrm t}}
\newcommand{\rad}{\operatorname{rad}}
\newcommand{\diam}{\operatorname{diam}}
\newcommand{\ecc}{\operatorname{ecc}}
\newcommand{\dist}{\operatorname{dist}}
\newcommand{\av}{\overline d_C}
\title{Distances between central vertices in graphs with small total domination number}
\author{AI-generated research manuscript}
\date{26 November 2026}
\begin{document}
\maketitle
\begin{abstract}
For a connected graph let $\av$ be the mean distance between distinct central
vertices, with value zero for a singleton center. We prove
$\av<\gt(G)$, the floor inequality posed in Graffiti.pc conjecture 268.
When $\gt(G)=2$ and the center has $m\ge2$ vertices, we determine the full set
of attainable means. Its largest member is $2-2/m$. Every value is realized
on $m+2$ vertices, the smallest possible order for a realization with a proper
center. We also characterize the induced center graphs in this class by a
construction preserving a specified dominating edge. When both ends of
that edge have pendant neighbors, four additional vertices are necessary
and sufficient. This gives an infinite family for which preserving the
edge increases the minimum from two to four.
\end{abstract}

\section{The question}
All graphs are finite, simple and undirected. Distances between centers are
taken in the ambient graph. Write $C=C(G)$ for the vertices of minimum
eccentricity, $m=|C|$, and $M=\binom m2$.
For $m\ge2$ put
\[
   \av(G)=\frac{1}{M}\sum_{\{x,y\}\subseteq C,\ x\ne y}d_G(x,y).
\]
For $m=1$ set $\av=0$. Averaging over ordered distinct pairs gives the same
number. A total dominating set $D$ supplies a neighbor in $D$ to \emph{every}
vertex, including the vertices of $D$. Write $\gt$ for its minimum size.
In a connected graph this number is finite exactly when there are at least
two vertices. Equivalently our strict inequality holds for every total
dominating set of every nonempty connected graph. This formulation treats the
one-vertex graph vacuously. With the extended-value convention it has
$\gt(K_1)=+\infty$.

DeLaVi\~na's catalogue \cite{wall} asks whether
$\lfloor1+\av\rfloor\le\gt$ (no.\ 268) and whether the ceiling suffices in
graphs containing no 4-cycle (no.\ 269). The latter forbids a 4-cycle even if
it has a chord. DeLaVi\~na, Liu, Pepper, Waller and West \cite{dlpww}
already proved $\rad(G)\le\gt(G)$. Their proof uses a total-domination-preserving
spanning tree, and the proof of their Theorem~5 uses adjacent central vertices
of a bicentral spanning tree. We give a component-sensitive derivation and use
the resulting strict inequality to address no.\ 268. Our extremal and
realization questions concern the internal distances and edges of $C$.

\section{A central edge in the equality case}
\begin{lem}\label{lem:spine}
Let $G$ be a connected graph, let $D$ be a total dominating set of size $k$, and let $c$ be the
number of components of $G[D]$. Then $G$ has a spanning tree of diameter at
most $k+2c-1$. In particular $2c\le k$.
\end{lem}
\begin{proof}
Every component of $G[D]$ has at least two vertices. Assign each vertex
outside $D$ to a component containing one of its neighbors and assign each
vertex of $D$ to its own component. Contract these parts. Each part is
connected, so the quotient is connected. For every edge in a spanning tree
of the quotient, choose one crossing edge $xy$ of $G$. Its ends lie either
in their assigned components or next to them. Adding the ends of these
$c-1$ edges to $D$ produces a connected dominating set $K$ with
$|K|\le k+2(c-1)$. Take a tree on $K$ using edges of $G[K]$ and attach
each remaining vertex as a leaf at one of its neighbors in $K$. A simple
path in the resulting tree contains at most $|K|$ vertices of $K$ and at most
two outside it. Its length is at most $|K|+1$.
\end{proof}

\begin{thm}\label{thm:268}
Every connected graph admitting a total dominating set satisfies
$\av(G)<\gt(G)$ and hence $\lfloor1+\av(G)\rfloor\le\gt(G)$.
If $k=\gt(G)$ is odd, then $\av(G)\le k-1$.
\end{thm}
\begin{proof}
For a minimum $D$ apply Lemma~\ref{lem:spine} to get a spanning tree $T$ of
diameter at most $2k-1$, or at most $2k-2$ when $k$ is odd.
The middle vertex or vertices of a longest tree path have eccentricity
$\lceil\diam(T)/2\rceil$. Indeed, if a vertex branches off the path at
position $i$ on a path of length $d$, its distance to the path is at most
$\min(i,d-i)$ by maximality. This proves the asserted upper bound for a
midpoint, and the two endpoints give the matching lower bound.
Distances can only decrease on passage from $T$ to $G$. Consequently
$r=\rad(G)\le k$, with $r\le k-1$ for odd $k$.

Every distance averaged in $\av$ is at most $r$. If $r<k$, the strict bound follows, including a singleton center under our convention. It remains to consider
$r=k$. Then $T$ must have diameter $2k-1$. Its two adjacent middle
vertices $u,v$ both have eccentricity $k$ in $T$ and hence in $G$,
as $k$ is already the minimum eccentricity in $G$. Thus $u,v\in C$.
Their mutual distance $1$ is strictly less than $k$ since every total
dominating set has at least two members. The finite average is strictly
less than $k$.
Finally $\lfloor 1+a\rfloor\le k$ is equivalent to $a<k$ for integer $k$.
\end{proof}

\section{A complete spectrum for a dominating edge}
A \emph{dominating edge} $uv$ has $N(u)\cup N(v)=V(G)$, using open
neighborhoods. Its endpoints totally dominate the graph. Thus its existence
is equivalent to $\gt=2$.

\begin{thm}\label{thm:spectrum}
Fix $m\ge2$. Among connected graphs with $\gt=2$ and exactly $m$ central
vertices, the set of mean center distances is exactly
\[
  \left\{\,2-\frac{j}{M}\ \middle|\ j=m-1,m,\ldots,M\,\right\}.
\]
Its maximum is $2-2/m$. Equality at the maximum holds exactly when
$G[C]$ is a tree. Each value is achieved by a graph with precisely two
noncentral vertices. This is best possible among graphs with a proper center.
\end{thm}
\begin{proof}
Let $uv$ be a dominating edge. Both endpoints have eccentricity at most
two. If the radius is one, every central vertex is universal and $G[C]$
is complete. If the radius is two, $u,v$ are central and give a dominating
edge of $G[C]$. In both cases $G[C]$ is connected, so its edge count $e$
is at least $m-1$. The distances of distinct centers are either one or two
and exactly $e$ unordered pairs have distance one. Thus
$M\av=e+2(M-e)=2M-e$. Equality at $e=m-1$ characterizes trees in a
connected graph.

For sharpness start with $K_{2,m-1}$, with part $\{u,q\}$ and other part
$S$. Attach one pendant vertex $p$ to $u$. For any $t$ from zero through
$\binom{m-1}{2}$ add any $t$ edges inside $S$. An edge from $u$ to any
$s\in S$ dominates all vertices. The vertex $p$ is at distance three from
$q$ since its unique neighbor $u$ misses $q$, and $p,u,s,q$ is a path.
Every $s\in S$ is within distance two of every vertex, and is at distance
two from $p$. The vertex $u$ has eccentricity two, witnessed by $q$.
Therefore $C=\{u\}\cup S$, and $e=m-1+t$. This realizes every claimed
integer $j$, including both endpoints; it also covers $m=2$.
Finally a vertex of maximum eccentricity in a graph with a proper center
has a farthest vertex. Both endpoints of this diametrical pair are
noncentral, so at least two additional vertices are necessary.
\end{proof}

For example, a five-vertex center gives exactly $1,11/10,6/5,13/10,7/5,
3/2,8/5$. The largest value comes from $K_{2,4}$ with a pendant vertex.
As $m$ grows, the upper endpoint tends to two. Thus no positive constant
can be subtracted uniformly from $\gt$ in Theorem~\ref{thm:268}.

\begin{thm}\label{thm:real}
A graph $H$ with at least two vertices is the induced center of a connected
graph of total domination number two if and only if $H$ has a dominating
edge. In the reverse direction any specified dominating edge of $H$ remains
dominating in a realization with four additional vertices, radius two and
diameter three.
\end{thm}
\begin{proof}
Necessity follows from the first paragraph of Theorem~\ref{thm:spectrum}.
For sufficiency let $uv$ be the specified edge. Preserve $H$ as an induced
subgraph and add four vertices $P,Q,A,B$. Join $P$ to every old vertex
except $v$, and $Q$ to every old vertex except $u$. Give $A$ exactly
the two neighbors $v,Q$ and $B$ exactly the two neighbors $u,P$.
These are all the additions.

The edge $uv$ still dominates, so all graph distances are at most three.
Any two old vertices other than the pair $u,v$ share $P$ or $Q$ if they
are nonadjacent. Every old vertex lies within two of both $P$ and $Q$,
using $uv$ for the two missing adjacencies. An old vertex other than
$u,v$ sees $A$ through $Q$ and $B$ through $P$; the endpoints use $uv$
where needed. No old vertex is universal, since each misses $A$ or $B$.
All therefore have eccentricity two. The pair $P,A$ is nonadjacent and
has no common neighbor, as is the pair $Q,B$. These four vertices have
eccentricity three. This proves the assertion, including its order bound.
\end{proof}
Buckley, Miller and Slater proved that every graph has a center realization
with at most four additional vertices \cite{bms}. The constraint here is
preservation of a specified total dominating pair. The unconstrained
appendage bound is classical.

\begin{thm}\label{thm:obstruction}
Let $H$ have a dominating edge $uv$ and let distinct leaves $x,y$, outside
$\{u,v\}$, be attached to $u,v$, respectively. Among all realizations of
$H$ as the induced center that preserve $uv$ as a dominating edge, the
minimum number of additional vertices is four.
\end{thm}
\begin{proof}
Consider any such realization $G$. Its radius is at most two. The center
contains the nonadjacent vertices $x,y$, so its radius is two. Every
noncentral vertex therefore has a vertex at distance greater than two,
and that vertex is also noncentral. Since $x,y$ are central, they have a
common neighbor $p$. It is noncentral, because their only neighbors in the
center are the distinct vertices $u,v$. By interchanging the ends if
necessary, suppose that $p$ is adjacent to $u$.

Choose $q$ at distance greater than two from $p$. It is adjacent to
neither $x$ nor $u$, both of which are adjacent to $p$. Since $x$ is
central, it has a common neighbor $t$ with $q$. The vertex $t$ is
noncentral, for $u$ is the only central neighbor of $x$ and misses $q$.
Also $t\ne p$ since $p$ misses $q$. Finally choose $z$ at distance greater
than two from $t$. The path $t,x,p$ and the edge $tq$ give $z\ne p,q$.
Thus $p,q,t,z$ are four distinct noncentral vertices. The matching
construction is Theorem~\ref{thm:real}.
\end{proof}

In particular let $H$ be a double star with at least one leaf on each end
of its central edge. Its ordinary centralizing number is two, as also
follows from the tree classification of Buckley, Miller and Slater
\cite{bms,behrens}. To see this directly, add one vertex adjacent to all
leaves and another adjacent to the two original internal vertices. These
two new vertices are at distance three. Every old vertex has eccentricity
two. At least two additions are necessary since $H$ is not self-centered
and a diametrical pair is outside any proper center. Theorem~\ref{thm:obstruction}
shows that requiring the original central edge to dominate increases the
minimum from two to four throughout this infinite family.

\section{The ceiling variant}
\begin{cor}\label{cor:269}
The ceiling inequality in no.\ 269 holds whenever $\gt$ is odd, or when
$\gt=2$ and the graph has no 4-cycle.
\end{cor}
\begin{proof}
The odd case follows from Theorem~\ref{thm:268}. In the second case choose a
dominating edge $uv$. If there is a universal vertex then $C$ is a clique.
Otherwise let $U,V$ consist of the vertices adjacent exclusively to $u,v$,
respectively, other than the endpoints. Both $U,V$ are nonempty.
There is at most one common neighbor of $u,v$. An edge from $U$ to $V$
gives a 4-cycle through $uv$, and an edge from the common neighbor to
either exclusive side also gives a 4-cycle. Thus every $U$--$V$ distance
is three: there is neither an edge nor a common neighbor between such
vertices. The only central vertices are $u,v$ and their common neighbor
if present, each within two of everyone. They form a clique. In all cases
$\av\le1$, giving the assertion.
\end{proof}
The even case $\gt\ge4$ of no.\ 269 is not settled here.

\section*{Research record}
The arguments and manuscript were generated with OpenAI Codex. The companion
archive contains exact-arithmetic experiments, a second enumerator, source and
priority notes, and the scope and execution status of the Lean artifacts.
The experiments count labelled graphs and do not enter the proofs.
A literature search located the older radius and center-realization results
cited above. Priority for the precise distance spectrum and constrained
realization and sharp appendage results remains open to bibliographic and
expert review.
\begin{thebibliography}{9}
\bibitem{wall} E. DeLaVi\~na, \emph{Written on the Wall II: Conjectures of
Graffiti.pc}, nos.\ 268--269, 23 February 2007. Catalogue updated 19 November
2026. \url{https://www.uhd.edu/academics/sciences/mathematics-statistics/wowii/index.html}.
\bibitem{dlpww} E. DeLaVi\~na, Q. Liu, R. Pepper, B. Waller, and D. B. West,
Some conjectures of Graffiti.pc on total domination, \emph{Congressus
Numerantium} \textbf{185} (2007), 81--95.
\url{https://www.uhd.edu/documents/academics/sciences/total-dom-2007.pdf}.
\bibitem{bms} F. Buckley, Z. Miller, and P. J. Slater, On graphs containing a
given graph as center, \emph{Journal of Graph Theory} \textbf{5} (1981),
427--434. \url{https://doi.org/10.1002/jgt.3190050413}.
\bibitem{behrens} S. L. Behrens, \emph{Graph Centers, Hypergraph Degree
Sequences, and Induced-saturation}, Ph.D. thesis, University of
Nebraska--Lincoln, 2015, Chapter~2.
\url{https://digitalcommons.unl.edu/mathstudent/63/}.
\end{thebibliography}
\end{document}

